\chapter{多模态模型}
\begin{quote}\small
\textit{本章摘要。}本章从图文检索、零样本分类和图像问答的输入输出出发，深入解释三个代表性模型。CLIP 用双塔编码器和对称对比损失学习图文匹配。BLIP-2 用查询变换器从视觉特征中提取信息，连接冻结的语言模型进行条件生成。LLaVA 将视觉表示接入语言上下文，通过视觉指令微调学习按问题回答。各模型的结构、训练目标和推理过程分别展开，最后讨论它们在工业迁移中怎样验证。
\end{quote}
视觉章节说明了如何从图像预测类别、框和像素边界。上一章说明了语言模型如何根据前文生成后续词元。现在给出一张巡检照片和问题“表面有什么异常”，还需要让语言模型取得照片中的信息。给出一句“有细长裂纹的工件”，则可能只需要从图库中找出相关照片。这两项任务都使用图像和文本，但输出形式与训练数据不同。

本章以图像与文本为主。先解释怎样比较二者，再解释怎样让图像成为语言生成的条件，最后解释怎样根据不同指令组织回答。阅读时需要用到内积、矩阵乘法、softmax、交叉熵和反向传播。视觉块表示与语言词元沿用前两章的定义。向量默认是列向量，将多个表示排成矩阵时，每行存放一个向量的转置。

\section{先确定任务：寻找匹配还是生成回答}
模态（modality）是信息的表现形式，例如图像、文本、音频和传感器序列。多模态模型处理两种或更多形式的输入，或者在不同形式之间建立对应。输入形式增加后，仍要先确定输出是什么。

\paragraph{图文检索与零样本分类。}
图文检索可以输入一句描述，输出按相关程度排序的图像，也可以输入一张图像，输出相关文本。模型需要给候选组合打分。它不必逐词生成描述。零样本分类则把候选类别写成文本，例如“正常表面的照片”“有裂纹的表面的照片”“有划痕的表面的照片”，再选出与图像最匹配的类别描述。这要求图像表示与文本表示能在同一空间比较。

\paragraph{图像描述与视觉问答。}
图像描述输入照片，输出一段概括内容的文字。视觉问答还输入问题，例如“裂纹大致位于左侧还是右侧”，输出针对该问题的回答。同一照片面对不同问题，需要使用不同信息。训练数据因此需要目标描述或问题与示范回答，仅有图文配对索引不足以规定怎样回答。

\paragraph{视觉指令跟随。}
用户还可能要求“只列出可见异常，不判断成因”或“按部件、现象、待确认信息三个字段回答”。模型既要读取图像，也要按指令选择内容和格式。视觉指令微调用这样的示范教模型响应请求。模型能回答位置问题，不等于具有经过验证的检测或分割能力。定位仍需按视觉章节的框或像素标注验收。

接下来选择 CLIP、BLIP-2 和 LLaVA 作为里程碑模型。这是一条解释能力差异的顺序，各模型拥有各自的训练方案。CLIP 学习匹配分数，BLIP-2 建立视觉条件生成路径，LLaVA 进一步学习视觉指令响应。本章中的图文“对齐”指表示与配对关系的对应，与上一章通过指令或偏好调整回答行为的含义需要区分。

% auxiliary-resource-guide
本章末的辅助学习材料按三个模型组织，可在完成相应算例后对照阅读。

\section{CLIP：把图像和文本变成可比较的表示}
\subsection{双塔为什么适合检索}
图文对比预训练（Contrastive Language--Image Pre-training，CLIP）用大量配对图像和文本学习共同表示空间\cite{radford2021clip}。训练文本可以是对图像内容的自然语言描述，无需先把所有内容限制成一套固定类别。共同空间要求两路输出可以用同一种内积比较，不要求像素和文字具有相同的原始格式。

设一个训练批次有$N$对样本$(v_i,t_i)$，其中$v_i$是第$i$幅图像，$t_i$是其配对文本，$i=1,\ldots,N$。视觉编码器$f_v$和文本编码器$f_t$各自处理输入，连同投影层输出$d$维表示：
\begin{equation}
 u_i^v=f_v(v_i)\in\R^d,\qquad
 u_i^t=f_t(t_i)\in\R^d.
\end{equation}
这里$d$是共同空间的宽度。视觉分支可以用卷积网络或视觉 Transformer，文本分支处理词元序列。两路编码器在计算单个表示时不读取另一模态，这种结构称为双塔。

若直接用未归一化向量的内积，向量长度也会影响分数。CLIP 将非零输出除以欧氏长度，得到单位向量：
\begin{equation}
 z_i^v=\frac{u_i^v}{\|u_i^v\|_2},\qquad
 z_i^t=\frac{u_i^t}{\|u_i^t\|_2}.
 \label{eq:multimodal-clip-normalization}
\end{equation}
$z_i^v,z_i^t\in\R^d$均有单位范数。式\eqref{eq:multimodal-clip-normalization}要求原向量非零。归一化后，内积等于两向量夹角的余弦，取值在$[-1,1]$之间。方向越接近，分数越高。模型通过训练学习这些方向，因此相近方向是否对应相关内容，要由配对数据和验证结果判断。

将所有图像表示按行排成$Z_v\in\R^{N\times d}$，文本表示排成$Z_t\in\R^{N\times d}$。全批次的得分为
\begin{equation}
 S=\frac{Z_vZ_t^{\mathsf T}}{\tau}\in\R^{N\times N},\qquad
 S_{ij}=\frac{(z_i^v)^{\mathsf T}z_j^t}{\tau}.
 \label{eq:multimodal-clip-scores}
\end{equation}
$\tau>0$是温度参数，控制得分差异进入 softmax 后的尺度，原始 CLIP 在训练中学习其对应尺度。$S$的行编号是图像，列编号是文本。$N$对观测产生$N^2$个组合，只有对角位置是按当前索引给定的配对。这里$S$已经除以温度，后续损失不再重复除以$\tau$。

图\ref{fig:multimodal-dual-encoder}以两对样本展示双塔如何形成得分表。两个编码器先独立产生单位表示，随后计算全部内积。方框标出的是训练目标位置，行归一化和列归一化分别解决两个检索方向的问题。
\begin{figure}[htbp]
\centering
\begin{tikzpicture}[x=1cm,y=1cm,every node/.style={font=\small}]
 \node[idi input,text width=1.3cm] (img) at (0,1.4) {$v_1,v_2$\\两幅图像};
 \node[idi input,text width=1.3cm] (txt) at (0,-0.3) {$t_1,t_2$\\配对描述};
 \node[idi process,text width=1.6cm] (ev) at (2.2,1.4) {$f_v$与归一化\\$Z_v$};
 \node[idi process,text width=1.6cm] (et) at (2.2,-0.3) {$f_t$与归一化\\$Z_t$};
 \draw[idi arrow] (img)--(ev);
 \draw[idi arrow] (txt)--(et);
 \node at (5.35,2.05) {$S=Z_vZ_t^{\mathsf T}/\tau$};
 \node at (5.0,1.5) {$t_1$}; \node at (6.1,1.5) {$t_2$};
 \node at (4.25,0.95) {$v_1$}; \node at (4.25,-0.05) {$v_2$};
 \draw (4.5,-0.55) rectangle (6.65,1.45);
 \draw (5.55,-0.55)--(5.55,1.45);
 \draw (4.5,0.45)--(6.65,0.45);
 \draw[thick] (4.63,0.59) rectangle (5.42,1.31);
 \draw[thick] (5.68,-0.41) rectangle (6.52,0.31);
 \node at (5.0,0.95) {$1$}; \node at (6.1,0.95) {$0$};
 \node at (5.0,-0.05) {$0$}; \node at (6.1,-0.05) {$1$};
 \draw[idi arrow] (ev.east)--(4.5,1.1);
 \draw[idi arrow] (et.east)--(4.5,-0.3);
 \draw[idi arrow] (6.8,0.95)--(7.7,0.95);
 \node[anchor=west,font=\scriptsize] at (7.0,0.55) {按行选文本};
 \draw[idi arrow] (6.1,-0.7)--(6.1,-1.35);
 \node[anchor=east,font=\scriptsize] at (5.85,-1.05) {按列选图像};
\end{tikzpicture}
\caption{两对样本的 CLIP 计算。箭头表示表示计算或归一化方向，方框标出给定配对。表中数值采用正文的二维单位向量与$\tau=1$，不代表实际模型的实验结果。}
\label{fig:multimodal-dual-encoder}
\end{figure}

\subsection{对称对比损失怎样使用配对索引}
固定图像$i$，可以把第$i$行的$N$个文本得分变成候选概率：
\begin{equation}
 P_{v\to t}(j\mid i)=\frac{e^{S_{ij}}}{\sum_{k=1}^N e^{S_{ik}}}.
\end{equation}
$j,k$都是文本候选编号。softmax 沿行计算，各概率非负且和为一。当前批次指定文本$i$为目标，因此用负对数概率惩罚模型没有选中它的程度。对所有图像取平均，得到图到文损失。

反过来，固定文本$i$，模型需要在所有图像中选出图像$i$。此时沿第$i$列归一化：
\begin{equation}
 P_{t\to v}(j\mid i)=\frac{e^{S_{ji}}}{\sum_{k=1}^N e^{S_{ki}}}.
\end{equation}
这里$j,k$变成图像候选编号。两方向使用相同的得分表，但候选轴不同。分别写出两项，再平均，就得到 CLIP 的对称对比损失：
\begin{align}
 \mathcal L_{v\to t}&=-\frac1N\sum_{i=1}^N
 \log\frac{e^{S_{ii}}}{\sum_{j=1}^N e^{S_{ij}}},\\
 \mathcal L_{t\to v}&=-\frac1N\sum_{i=1}^N
 \log\frac{e^{S_{ii}}}{\sum_{j=1}^N e^{S_{ji}}},\\
 \mathcal L_{\rm CLIP}&=\frac12
 (\mathcal L_{v\to t}+\mathcal L_{t\to v}).
 \label{eq:multimodal-clip-loss}
\end{align}
这些量是当前训练批次的经验损失，使用自然对数。目标索引来自图文配对，正样本指配对组合，负样本指被用作不匹配候选的其他组合，与“正常”“故障”的业务类别无关。图文必须同步重排，或者同步更新目标索引，否则对角目标就不再正确。

为了看清参数为何会变化，固定图像$i$，记单行损失为$\ell_i=-\log P_{v\to t}(i\mid i)$。展开对数并对一个得分求导：
\begin{align}
 \ell_i&=-S_{ii}+\log\sum_{k=1}^N e^{S_{ik}},\\
 \frac{\partial\ell_i}{\partial S_{ij}}
 &=P_{v\to t}(j\mid i)-\mathbf1_{j=i}.
\end{align}
$\mathbf1_{j=i}$在$j=i$时为一，其他情况为零。若得分可独立调整，梯度下降会提高目标得分、降低其余得分。实际得分共享编码器参数，一次更新汇总所有行列的梯度，不能保证每个组合都单独按这个方向变化。

\subsection{两对样本：从内积算到损失}
取$N=2,d=2$，设两幅图像的单位表示依次为$(1,0)^{\mathsf T}$、$(0,1)^{\mathsf T}$，两段配对文本也取这两个向量。取$\tau=1$，则
\begin{equation}
 Z_v=Z_t=\begin{bmatrix}1&0\\0&1\end{bmatrix},\qquad
 S=\begin{bmatrix}1&0\\0&1\end{bmatrix}.
\end{equation}
四个组合中，配对内积为一，交叉组合内积为零。对第一幅图像，两个候选文本的概率为
\begin{equation}
 P_{v\to t}(1\mid1)=\frac{e}{e+1}\approx0.731,\qquad
 P_{v\to t}(2\mid1)=\frac1{e+1}\approx0.269.
\end{equation}
第二行以及两列的正确配对概率相同。因此两方向平均损失都约为$-\log0.731=0.313$，对称损失也约为$0.313$。两个方向已经分开，softmax 仍不会在有限得分下输出概率一。

保持表示不变，仅把$\tau$改为$0.5$，对角得分变为二，正确配对概率变为$e^2/(e^2+1)\approx0.881$，损失约为$0.127$。损失下降来自尺度改变，图文方向并未改善。温度因而不能脱离表示和任务指标单独评价。固定得分差时，小温度使概率更集中，也放大相似度梯度的尺度。它与上一章生成时调整词表分布的解码温度作用于不同对象。

若$N=1$且没有额外负样本，分子分母完全相同，两项损失都恒为零，无法提供对比分数梯度。较大的批次增加候选，但候选未必真的不匹配。例如两幅图都显示裂纹，两段文本都写“有裂纹的表面”，非对角组合也可能合理。这类实际匹配却被当成负例的组合称为假负例（false negative）。若数据存在多重匹配，应检查配对规则，并考虑将多个合理目标共同计入监督。

\subsection{训练之后怎样检索与分类}
训练时，配对文本提供监督，损失通过编码器和投影层更新参数。推理时参数固定，不再需要正确配对索引。文本检索图像可以先缓存图库的单位图像向量，再编码输入描述，按内积排序。图像检索文本则缓存候选文本向量。双塔使这两种缓存都可行，更新编码器后需重新生成对应索引。

零样本分类把类别名称放入描述模板。设有$C$个互斥候选类别，类别$c$的描述为$t_c$，单位文本向量为$z_c^t$。新图像的单位向量为$z^v$，决策写为
\begin{equation}
 a_c=(z^v)^{\mathsf T}z_c^t,\qquad
 \hat y=\argmax_{c\in\{1,\ldots,C\}}a_c.
\end{equation}
$a_c\in[-1,1]$是类别相似度，$\hat y$是选出的类别编号。正温度缩放不改变排序。这里“零样本”表示在这种迁移设置下，没有使用目标任务的标注图像训练类别预测器，模型仍经历过图文预训练。若随后用目标域标签选择提示或调节阈值，应报告这部分监督，不能继续把整个流程称为严格零样本。

对类别得分做 softmax，可以得到相对于这$C$类的归一化数值，但它不自动等于类别正确的校准概率。照片也可能含有候选之外的缺陷，或者同时出现裂纹和划痕。前者需要拒识检验，后者需要多标签任务定义。描述模板与候选集合都会改变结果。

CLIP 将整幅图像和整段文本各汇成一个向量，适合大规模匹配。整图相似度难以直接表达“左边的零件挡住右边的裂纹”这样的细粒度关系，也不提供逐词生成路径。要解释照片或回答问题，需要继续把视觉信息送入语言模型。

\section{BLIP-2：用少量查询连接视觉与语言}
\subsection{为什么需要一个可训练的连接模块}
已有视觉编码器能产生区域特征，已有语言大模型能生成文字，但两者的表示宽度和训练任务不同。把视觉数表直接当成文本嵌入，没有理由让语言模型理解。若将两个大模型从头联合训练，更新参数和训练资源又会明显增加。

第二代视觉语言预训练模型 BLIP-2（Bootstrapping Language-Image Pre-training，第二代）用一个较小的查询变换器（Querying Transformer，Q-Former）连接冻结的视觉编码器和语言模型\cite{li2023blip2}。冻结指训练时不更新指定权重。Q-Former 先学习从图像中提取与文本有关的信息，再学习把这些信息作为语言模型可以利用的条件输入。

设视觉编码器输出$H_v\in\R^{n_v\times d_v}$，其中$n_v$是视觉位置数，$d_v$是每个位置的特征宽度。沿用视觉章节的理解，这些位置可对应图像块，但特征已受编码器上下文影响，并非每行只含一个独立局部区域的信息。

Q-Former 使用$m$个可学习查询，初始查询矩阵为$Q_0\in\R^{m\times d_q}$。$d_q$是查询宽度，$m$由结构设定。查询是训练得到的向量，不是$m$句人工问题，也没有预先规定每个查询负责哪种物体。原始 BLIP-2 使用固定数量的查询，下面取较小数值说明计算。

\subsection{查询通过交叉注意力读取哪些信息}
查询先在自身之间进行自注意力交互，再通过交叉注意力（cross-attention）读取视觉特征。与视觉章节的自注意力相比，这里查询和键值来自不同的序列。令某层查询状态为$H_q\in\R^{m\times d_q}$。单头交叉读取可写为
\begin{align}
 Q&=H_qW_Q\in\R^{m\times d_k},\\
 K&=H_vW_K\in\R^{n_v\times d_k},\\
 V&=H_vW_V\in\R^{n_v\times d_a},\\
 A&=\softmax\left(\frac{QK^{\mathsf T}}{\sqrt{d_k}}\right),
 \qquad O=AV.
 \label{eq:multimodal-query-attention}
\end{align}
$W_Q\in\R^{d_q\times d_k}$、$W_K\in\R^{d_v\times d_k}$、$W_V\in\R^{d_v\times d_a}$为可训练投影。$d_k$是查询与键的共同宽度，$d_a$是值的宽度。$A\in\R^{m\times n_v}$逐行沿视觉位置归一化，$O\in\R^{m\times d_a}$保留查询行数。$\sqrt{d_k}$用于控制内积尺度，在各维独立且单位方差的简化条件下，未缩放内积的方差随$d_k$增长。

第$a$个查询读取结果为$O_{a,:}=\sum_{b=1}^{n_v}A_{ab}V_{b,:}$。$a$是查询编号，$b$是视觉位置编号，$A_{ab}\geq0$且$\sum_bA_{ab}=1$。不同查询可以从相同视觉特征读取不同组合。实际 Q-Former 使用多头注意力、输出投影、残差和前馈层，式\eqref{eq:multimodal-query-attention}只展开其中一次读取。多头结果需投影回$d_q$宽度，才能与原查询状态做残差相加。

\paragraph{两个查询读取三个位置的手算。}
取$m=2,n_v=3,d_k=1,d_a=1$。用以下查询和键构造明确的示意得分：
\begin{equation}
 Q=\begin{bmatrix}1\\-1\end{bmatrix},\qquad
 K=\begin{bmatrix}\log2\\0\\0\end{bmatrix},\qquad
 V=\begin{bmatrix}2\\4\\10\end{bmatrix}.
\end{equation}
$V$的三个标量是教学特征值，没有像素或毫米单位。第一行得分为$(\log2,0,0)$，softmax 权重为$(1/2,1/4,1/4)$。第二行得分为$(-\log2,0,0)$，权重为$(1/5,2/5,2/5)$。于是
\begin{equation}
 A=\begin{bmatrix}1/2&1/4&1/4\\1/5&2/5&2/5\end{bmatrix},
 \qquad O=\begin{bmatrix}4.5\\6\end{bmatrix}.
\end{equation}
第一个结果为$2/2+4/4+10/4=4.5$，第二个为$2/5+8/5+20/5=6$。输出有两行，与查询数一致，不能当成三个视觉位置的标签。若另一组查询与键产生权重$(0.1,0.2,0.7)$，读取结果就会变为$0.1\times2+0.2\times4+0.7\times10=8$。这个变化说明读取结果依赖查询与视觉特征共同形成的权重。

图\ref{fig:multimodal-query-readout}画出本次计算的两个输入轴。每个查询连接三个值，边旁数字表示读取权重。它们是模型计算中的加权系数，不提供真实缺陷位置或因果解释。
\begin{figure}[htbp]
\centering
\begin{tikzpicture}[x=1cm,y=1cm]
 \node[idi input,text width=1.6cm] (v1) at (0,2.4) {位置1\\值$2$};
 \node[idi input,text width=1.6cm] (v2) at (0,0) {位置2\\值$4$};
 \node[idi input,text width=1.6cm] (v3) at (0,-2.4) {位置3\\值$10$};
 \node[idi state,text width=1.7cm] (q1) at (4.6,1.6) {查询1读取\\$O_1=4.5$};
 \node[idi state,text width=1.7cm] (q2) at (4.6,-1.6) {查询2读取\\$O_2=6$};
 \draw[idi arrow] (v1.east)--node[pos=0.67,above,font=\scriptsize] {$1/2$} (q1.west);
 \draw[idi arrow] (v2.east)--node[pos=0.7,above,font=\scriptsize] {$1/4$} (q1.west);
 \draw[idi arrow] (v3.east)--node[pos=0.7,above,font=\scriptsize] {$1/4$} (q1.west);
 \draw[idi arrow,dashed] (v1.east)--node[pos=0.35,below,font=\scriptsize] {$1/5$} (q2.west);
 \draw[idi arrow,dashed] (v2.east)--node[pos=0.68,below,font=\scriptsize] {$2/5$} (q2.west);
 \draw[idi arrow,dashed] (v3.east)--node[pos=0.67,below,font=\scriptsize] {$2/5$} (q2.west);
 \node[anchor=west,font=\small,align=left] at (6.05,0) {三个位置\\压成两个输出\\实线：查询1\\虚线：查询2};
\end{tikzpicture}
\caption{单头交叉注意力的数值读取。箭头表示值信息进入查询输出，实线和虚线区分两个查询的读取路径。权重由正文给定的查询与键计算，输出行数由查询数决定。}
\label{fig:multimodal-query-readout}
\end{figure}

\subsection{两阶段训练分别解决什么问题}
\paragraph{第一阶段：从图像提取与文本有关的表示。}
这一阶段冻结视觉编码器，训练 Q-Former，还没有接入冻结的语言大模型。Q-Former 除查询路径外还有文本路径，两者共享部分 Transformer 参数，利用不同注意力掩码定义任务。原始方法共同使用图文对比、图文匹配和以图像为条件的文本生成三个目标。

图文对比要求查询输出与配对文本表示能匹配。原始方法比较文本表示与多个查询表示，并使用其中最大的相似度形成图文对比得分。此任务限制查询与文本路径直接互读，避免仅复制配对文本就得到好分数。图文匹配则让查询和文本交互，判断这组图文是否匹配，学习更细的联合关系。图像条件文本生成让文本位置读取查询输出，同时只能读取此前文本，以下一词元为目标。它要求查询携带有助于描述图像的信息。

三项目标分别约束独立表示的可比较性、联合图文的匹配判断和视觉条件下的文本预测。第一阶段的生成由 Q-Former 的文本路径承担，不能与第二阶段接入语言大模型的生成混为同一步。不同掩码改变了哪些信息可见，监督目标仍由配对关系或真实文本确定。

\paragraph{第二阶段：让冻结语言模型利用视觉条件。}
保留冻结视觉编码器，接入冻结语言模型，训练 Q-Former 和连接投影。记 Q-Former 的最终查询输出为$H_Q\in\R^{m\times d_q}$。线性投影得到
\begin{equation}
 B=H_QW_P\in\R^{m\times d_{\rm lm}},\qquad
 W_P\in\R^{d_q\times d_{\rm lm}}.
 \label{eq:multimodal-blip-prefix}
\end{equation}
$d_{\rm lm}$是语言模型嵌入宽度，$B$包含$m$个连续视觉条件向量。这里省略可选偏置。这些向量常称为软视觉前缀，它们占据输入位置，却不需要先变成词表中的离散编号。BLIP-2 可连接不同语言模型结构，下面选解码器式语言模型说明前缀计算。

给定图像$v$、文本提示$q$和目标回答$y_{1:T}$，条件概率和单样本训练损失为
\begin{align}
 p(y_{1:T}\mid v,q)
 &=\prod_{t=1}^T p(y_t\mid B(v),q,y_{<t}),\\
 \ell_{\rm gen}
 &=-\frac1T\sum_{t=1}^T
 \log p(y_t\mid B(v),q,y_{<t}).
 \label{eq:multimodal-generation-loss}
\end{align}
$B(v)$表示由该图像计算的视觉前缀，$q$是已知提示，$y_{<t}$是第$t$个目标词元之前的真实回答。$T\geq1$是计分回答词元数，可包含结束符，填充不计分。$y_t$取自语言模型词表。所有概率由带因果掩码的语言模型计算，回答位置可以读取视觉前缀、提示和此前回答。训练集目标由各样本的损失平均得到。

这沿用上一章的下一词元预测。图文对比损失在有限候选间选择配对，式\eqref{eq:multimodal-generation-loss}在每一步的词表中预测目标词元。二者的候选、标签和输出都不同。生成损失降低只说明示范文字的预测改善，不直接证明每句话得到图像支持。

\paragraph{冻结权重时，梯度在哪里停止。}
冻结语言模型参数$\theta_{\rm lm}$后，仍要计算它对视觉前缀的导数。令可训练连接部分的参数为$\phi$，链式法则为
\begin{equation}
 \nabla_\phi\ell_{\rm gen}
 =\left(\frac{\partial B}{\partial\phi}\right)^{\mathsf T}
 \nabla_B\ell_{\rm gen}.
\end{equation}
此处将$B$展平理解为向量，雅可比矩阵$\partial B/\partial\phi$把连接参数变化映射到视觉前缀变化。$\nabla_B\ell_{\rm gen}$来自穿过语言模型计算的反向传播。语言模型权重不执行更新，输入导数仍需保留。若在语言模型输入处切断梯度，连接模块就无法从回答损失学习。视觉编码器的权重也不更新，在本方案中其输出可以作为不求导的特征输入交给 Q-Former。

所以冻结减少的是需要更新的权重及相应优化器状态，冻结语言模型仍需前向计算，也通常仍需为输入梯度保存或重算内部激活。训练成本不能仅按可训练参数比例估计。

\subsection{查询压缩带来什么收益与限制}
取一个教学设置，视觉特征为$64\times128$，使用八个查询，查询宽度$d_q=32$，语言宽度$d_{\rm lm}=256$。Q-Former 输出$8\times32$，连接投影为$32\times256$，视觉前缀为$8\times256$。原有64个视觉位置压成八个语言输入位置。原始模型的配置不同，这些数值只用于追踪维度。

固定查询数可以控制送入语言模型的视觉序列长度，但 Q-Former 仍要读取64个视觉位置。压缩不能保证保留全部小缺陷、文字和空间关系。若视觉编码器在输入缩放时已丢失裂纹像素，连接模块没有原始证据可供恢复。增加查询数可能提高信息容量，也增加计算，需要在相同数据和预算下检验。

推理时不给真实回答，也不更新参数。先计算图像特征、查询输出和视觉前缀，再给提示，例如“描述可见表面”，逐词生成文字。若接入的是以描述训练为主的模型，面对复杂格式指令或多轮问题可能表现不足。视觉条件可用与指令响应熟练是两项不同要求，下一节用 LLaVA 解释后一项怎样训练。

\section{LLaVA：让视觉表示参与指令回答}
\subsection{从线性投影到语言上下文}
大型语言与视觉助手（Large Language and Vision Assistant，LLaVA）通过视觉指令微调，把预训练视觉编码器与语言模型连接起来\cite{liu2023llava}。原始 LLaVA 用线性投影连接二者，没有使用 Q-Former 的查询压缩结构。图像描述用于初步连接，随后用包含问题、指令和示范回答的数据训练对话响应。

设视觉编码器仍输出$H_v\in\R^{n_v\times d_v}$。原始线性连接可写为
\begin{equation}
 B_v=H_vW_C\in\R^{n_v\times d_{\rm lm}},\qquad
 W_C\in\R^{d_v\times d_{\rm lm}}.
 \label{eq:multimodal-llava-projection}
\end{equation}
$B_v$的每行是送入语言模型的一个视觉词元表示，这些连续向量没有普通文字的词元编号。投影改变宽度，保留$n_v$行。沿用$64\times128$的教学视觉特征，投影矩阵为$128\times256$，得到$64\times256$。与 BLIP-2 的八个查询前缀相比，这里直接接入64个位置。

后续常用的 LLaVA 版本将线性连接改为多层感知机（Multilayer Perceptron，MLP）。例如逐位置的两层连接可写为
\begin{equation}
 B_v=\sigma(H_vW_1)W_2,
 \label{eq:multimodal-llava-mlp}
\end{equation}
$W_1\in\R^{d_v\times d_h}$、$W_2\in\R^{d_h\times d_{\rm lm}}$，$d_h$为连接器隐藏宽度，$\sigma$是逐元素非线性激活，这里省略偏置。它能学习更灵活的逐位置变换，行数仍为$n_v$，不会自动产生查询压缩。该式用于说明后续变体，原始论文的连接器应按式\eqref{eq:multimodal-llava-projection}理解。

将视觉向量与提示的文字嵌入沿序列方向拼接。若提示$q$有$n_q$个词元，文字嵌入为$E(q)\in\R^{n_q\times d_{\rm lm}}$，一个简化上下文为
\begin{equation}
 X_{\rm ctx}=\begin{bmatrix}B_v\\E(q)\end{bmatrix}
 \in\R^{(n_v+n_q)\times d_{\rm lm}}.
\end{equation}
$X_{\rm ctx}$包含视觉位置和文字位置，宽度统一为$d_{\rm lm}$。实际对话还含角色与分隔标记，图像向量在图像占位位置插入。视觉与文本不按坐标相加，语言模型通过后续注意力读取它们。

图\ref{fig:multimodal-llava-context}用三个视觉位置和两个提示位置画出序列。生成第一个回答词元时，它可以读取全部已知视觉与提示位置。训练时即使将整段真实回答同时送入网络，因果掩码仍禁止该位置读取自己的目标和后续回答。
\begin{figure}[htbp]
\centering
\begin{tikzpicture}[x=1cm,y=1cm,every node/.style={font=\small}]
 \foreach \i in {1,2,3}{
  \node[idi input,minimum width=1.05cm,inner sep=3pt]
   (v\i) at ({1.35*(\i-1)},0) {$b_{\i}$};
 }
 \node[idi state,minimum width=1.05cm,inner sep=3pt] (q1) at (4.05,0) {$q_1$};
 \node[idi state,minimum width=1.05cm,inner sep=3pt] (q2) at (5.4,0) {$q_2$};
 \node[idi output,minimum width=1.05cm,inner sep=3pt] (y1) at (6.75,0) {$y_1$};
 \node[draw,dashed,align=center,minimum width=1.05cm,minimum height=0.72cm]
  (y2) at (8.1,0) {$y_2$};
 \node at (1.35,0.8) {投影后的视觉位置};
 \node at (4.72,0.8) {提示词元};
 \node at (7.4,0.8) {回答目标};
 \node[idi process,text width=3.0cm] (pred) at (4.05,-2.1)
  {根据已知上下文\\预测第一个回答词元};
 \foreach \s in {v1,v2,v3,q1,q2}{
  \draw[idi arrow] (\s.south)--(pred.north);
 }
 \draw[idi arrow] (pred.east)-| (y1.south);
 \node[font=\scriptsize,align=center] at (7.35,-2.0)
  {$y_1,y_2$不作为\\本步可见输入};
\end{tikzpicture}
\caption{视觉指令回答的上下文与目标。$b_1,b_2,b_3$表示三行投影视觉向量，$q_1,q_2$表示简化提示词元。箭头表示本步可读取的信息和预测输出。$y_2$的虚线框表示未来目标，图中省略角色标记和内部注意力层。}
\label{fig:multimodal-llava-context}
\end{figure}

\subsection{两阶段训练怎样改变回答能力}
\paragraph{第一阶段：学习视觉与语言的连接。}
原始 LLaVA 冻结视觉编码器和语言模型，只训练线性投影。样本提供图像和配对描述，以描述的下一词元为目标。模型必须通过连接后的视觉向量提高真实描述的概率，才能让视觉输出进入语言模型可利用的输入范围。这一阶段更新少量参数，但同样需要保留通过冻结语言模型回到投影的输入梯度。

\paragraph{第二阶段：学习按视觉指令回答。}
继续冻结视觉编码器，训练连接层和语言模型。样本包含图像、用户问题或指令以及示范回答，可含多轮对话。回答损失使用式\eqref{eq:multimodal-generation-loss}的形式，将视觉条件$B(v)$替换为$B_v(v)$。计分位置来自助手回答，问题、角色标记和视觉前缀作为条件。损失掩码决定哪些位置计分，因果注意力掩码决定哪些位置可读取，两者作用不同。

原始工作利用语言模型根据图像描述与框等已有文本标注辅助构造视觉指令数据。这样的示范扩大了问答形式，但内容质量仍受标注覆盖和生成误差影响。若示范把没有证据的成因写成肯定结论，下一词元训练也会鼓励类似表达。指令微调学习示范中的响应方式，不直接给出事实正确性保证。

这与 BLIP-2 的两阶段训练有明确差异。BLIP-2 第一阶段训练查询提取，第二阶段连接冻结语言模型。原始 LLaVA 第一阶段训练投影，第二阶段同时调整语言模型来适应视觉指令。后续使用参数高效微调或不同冻结策略时，应单独说明实际更新范围。

\subsection{同一张图为何能产生不同回答}
考虑一张教学照片，人工标注确认其右侧有一条细长裂纹。输入“描述这张照片”时，示范可以要求概括可见物体和异常。输入“只回答是否有裂纹”时，示范可以只输出“有”。输入“按现象和待确认信息回答，不推测成因”时，示范可以输出“现象：右侧可见细长裂纹。待确认：裂纹深度和形成原因”。这些是教学输出示例，不是模型实测结果。

三次输入使用相同视觉特征，不同提示改变语言上下文和后续条件概率。参数在推理期间固定。训练时真实回答前缀由数据提供，生成时则依次选词元、追加到上下文，再预测下一词元。若首句把划痕误写成裂纹，后续回答可能继续围绕这项错误组织解释。

幻觉（hallucination）在这里指生成了图像或已有证据不支持的内容。例如照片只显示表面，模型却给出“裂纹深度为两毫米”，既缺少深度观测，也缺少毫米尺度的依据。降低解码随机性不能补足这些输入。含有“不确定”示范可以训练相应行为，仍需验证模型是否在证据不足时实际使用它。

LLaVA 直接保留更多视觉位置，可能减少查询压缩造成的信息损失，但语言上下文更长。标准全局注意力对总序列长度的交互开销呈二次增长。保留位置也不保证细裂纹识别或准确定位，输入分辨率、特征提取、位置表示和监督粒度仍决定可用证据。注意力图不能直接替代真实框或像素标签。

\subsection{用相同维度比较 BLIP-2 与 LLaVA}
两种连接都从$H_v$出发，最终提供$d_{\rm lm}$宽的语言输入。BLIP-2 先让查询读取视觉特征，输出行数由$m$控制。原始 LLaVA 逐行线性投影，输出行数仍为$n_v$。前者在语言模型前压缩位置，后者将更多位置交给语言模型继续处理。

表\ref{tab:multimodal-interaction}按相同维度总结已经解释的差异。可训练参数少、上下文短和任务效果好是不同结论，需要分别检查。

\begin{table}[htbp]
\centering
\caption{BLIP-2 与原始 LLaVA 的连接和训练差异}
\label{tab:multimodal-interaction}
\small
\setlength{\tabcolsep}{4pt}
\begin{tabular}{@{}p{2.0cm}p{5.1cm}p{5.1cm}@{}}
\toprule
比较项 & BLIP-2 & 原始 LLaVA \\
\midrule
视觉连接 & 查询交叉读取，再投影 & 对每个视觉位置线性投影 \\
语言输入 & $m$个查询前缀 & $n_v$个投影视觉位置 \\
第一阶段 & 冻结视觉编码器，训练查询与文本路径 & 冻结视觉与语言模型，训练投影 \\
第二阶段 & 视觉与语言模型均冻结，训练查询和投影 & 视觉编码器冻结，训练投影和语言模型 \\
主要收益 & 控制视觉前缀长度，减少大模型权重更新 & 连接直接，通过指令数据调整响应方式 \\
主要限制 & 查询压缩可能丢失细节 & 视觉位置多，语言上下文成本较高 \\
\bottomrule
\end{tabular}
\end{table}

\section{典型应用怎样验证}
\subsection{先按输出形式选择评价对象}
图文检索应检查相关候选是否进入前$k$个结果。$k$是预先指定的返回数量，相关性由人工配对或任务标注确定。若一条描述对应多幅图像，应在评价中保留多重相关关系。训练批次的配对概率与全图库的检索效果不同，不能用低对比损失替代留出图库验证。

零样本分类应固定候选类别和提示模板，再在未用于选择提示的标注图像上统计准确率、各类召回和误报。细裂纹、反光、不同相机和类别之外的输入需分别检查。若使用目标域标签训练分类头，这就是有监督迁移，应与原来的零样本设置分开报告。

图像描述与视觉问答需要核对内容。可以把回答拆成可检验陈述，例如“有裂纹”“在右侧”“深度两毫米”，逐项检查图像或外部资料是否支持。还要统计应答覆盖：该回答是否遗漏了任务要求的可见现象。语言流畅度、示范词元损失和事实正确率分别评价不同对象。

视觉指令跟随还要检查格式和限制是否遵守。例如要求不推测成因时，即使类别判断正确，添加未经证实的故障成因也不满足任务。位置回答则另用框、像素标注或约定的位置类别验收。模型名称和注意力可视化不能替代这些标签。

\subsection{工业迁移：从巡检照片到报告草稿}
考虑一个范围有限的应用：检索相似巡检照片，筛查表面缺陷，再草拟可见现象报告。CLIP 可先按描述检索图库或比较候选缺陷类别。需要文字描述时，可以评估 BLIP-2。需要用户指定字段或多轮追问时，可以评估 LLaVA。它们承担不同输出，不要求为了流程完整而全部部署。

验证集与测试集应按设备、生产批次或事件隔离。同一工件的近重复照片不能跨越训练与测试，提示、阈值和微调方案也不能在测试集上挑选。迁移时比较已有视觉分类器、CLIP 零样本结果和目标域微调结果，保持数据划分与标注定义一致，再记录可训练参数、视觉位置数和生成延迟。这样才能分辨收益来自新增监督、结构差异还是更多计算。

外部文字也要明确来源。有效手册可以支持报告中的规程说明，模型仍需按上一章的检索方式确认版本和适用条件。维修结论可以成为事后监督，预测时尚未形成的结论不能放进输入。沿用原有时间例子：10:00诊断时，09:59:58拍摄但10:00:03才到达的照片不可用。10:30确认的类别可以作为历史样本标签，不能回填成10:00已知信息。

若额外接入振动或工艺曲线，需要为该联合任务提供对应数据和标签，图文预训练本身不建立传感器理解能力。假设一百个事件中图像分类答对八十个、振动分类答对七十五个，两个数不能推出融合会答对九十个。要逐事件比较错误是否互补，再与单模态和简单组合基线比较。照片缺失、模糊或两路信息冲突时还应分别测试，证据不足允许请求复检。

所有数值算例和问答示例用于解释机制，均不代表已完成的工业实验。报告草稿应保留图像来源与待确认信息，最终结论由独立证据和任务验收决定。

\section*{本章小结}
CLIP 用配对索引训练两个独立编码器，在单位表示的共同空间中形成图文得分。对称损失同时训练图到文和文到图选择，推理只需编码和排序。它适合检索与零样本分类，缺少直接生成回答的路径。

BLIP-2 让少量可学习查询通过交叉注意力读取视觉特征，再将查询输出投影为语言模型的条件。其两阶段训练先学习视觉语言表示，再连接冻结语言模型。冻结权重保留了输入梯度，查询压缩则引入信息容量与上下文成本的折中。

LLaVA 将视觉位置投影到语言嵌入宽度，通过描述数据学习连接，再通过视觉指令数据调整连接层和语言模型。原始线性投影与后续 MLP 需要区分。它可以根据问题组织回答，回答是否忠于图像、是否准确定位仍要独立验证。

本章的生成是在已知图像条件下产生文字。下一章转向图像、波形等数据的生成，研究怎样从随机变量或噪声出发构造新样本，并检查所学分布与真实数据的关系。

\section*{辅助学习材料}
\begingroup
\emergencystretch=3em
\begin{itemize}[leftmargin=*]
\item \textbf{CLIP 原始论文}\cite{radford2021clip}。对照双塔结构、批内对称损失和零样本分类，先重算本章两对样本，再查看论文中的真实配置与评价设置。
\item \textbf{BLIP-2 原始论文}\cite{li2023blip2}。重点读 Q-Former、第一阶段三个目标的注意力掩码，以及第二阶段连接不同语言模型的方式。将论文中实际的查询数与本章教学维度区分。
\item \textbf{Visual Instruction Tuning}\cite{liu2023llava}。对照原始线性连接、两阶段冻结范围和视觉指令数据构造。后续代码配置中的 MLP 不应反推成原始模型结构。
\item \textbf{Contrastive Representation Learning}（Lilian Weng，英文教程）。复习对比目标、温度与负样本选择，对照 CLIP 损失和假负例讨论。\href{https://lilianweng.github.io/posts/2021-05-31-contrastive/}{在线阅读}。
\item \textbf{Natural language image search with a Dual Encoder}（Keras，英文代码教程）。追踪双编码器、投影和检索计算。示例与 CLIP 的具体训练配置不同，运行前核对依赖版本。\href{https://keras.io/examples/vision/nl\_image\_search/}{在线阅读}。
\end{itemize}
\endgroup
